\scanpage{7}
\mspage{1}

Soit \(M\) une catégorie (``\unsure{modèles|modéles}''). On suppose

(i) Toute flèche de \(M\) est monomorphique

(ii) Tout \deleted{autom}\inserted{end} d'un élément de \(M\) est l'identité.

Ex : \(M\) une catégorie définie \unsure{par} un ensemble ordonné

Pour tout \(D\in\mathrm{Ob}\,M\), soit \(\widetilde D\) l'ensemble (ordonné par ``inclusion'') des sous objets de \(D\). On identifie \(D\) au plus grand élément de \(\widetilde D\)

Soient \(D',D\in\mathrm{Ob}\,M\), alors l'application
\begingroup\makeatletter\tagsleft@true\makeatother
\[
\begin{aligned}
\underset{M}{\mathrm{Hom}}(D',D)&\longrightarrow\widetilde D\\
u&\longmapsto u(D')
\end{aligned}
\tag*{(\(\ast\))}
\]
\endgroup
est injective par \unsure{[i]|(i)}(ii), son image est formée de l'ens \deleted{\unsure{de}}
\(\Phi_{D'}(D)\subset\widetilde D\)
des sous-objets \(\Delta\) de \(D\) qui sont \inserted{\(M\)-}isomorphes à \(D'\).

\deleted{Soit \(K\) un ens. ordonné}

Soit \(N\)\deleted{\(I\)} l'ens \inserted{préordonné} des classes d'isom. d'objets de \(M\). \inserted{Il est en fait \uline{ordonné}, grâce à (ii).}
Pour \(\unsure{\forall|\mathrm H}\ n\in\mathbb N\), un\deleted{\unsure{e|s}} objet \(D_n\) de \deleted{\(D\)}\(M\) de classe \(\overset{\ast}{n}\) est déterminé à isom unique près, \unsure{à|\(\therefore\)} cause de (ii), donc l'ens ordonné \(\widetilde D_n=\mathbb J_n\) est défini à isom unique près. On a donc \deleted{\unsure{plus}, une}\inserted{une} famille
\[
(\mathbb J_n)_{n\in\mathbb N}
\]
d'ens ordonnés, indexés par l'ensemble ordonné \(\mathbb N\).
% archive sheet number: 6
